\section{Evidence packing attribution and influence}\label{sec:C16:packing}
Retrieved content becomes useful only after it reaches the model in an interpretable form. Packing determines ordering, redundancy, source labels, and the separation between evidence and instruction. A strong ranker can still fail if the prompt truncates its best result or strips the qualification that made it relevant. The packed context should be treated as a versioned artifact that can be reconstructed for audit.

MemCalib studies a downstream distinction: a memory atom can be appropriately ignored, provide bounded support, or control a material conclusion~\cite{memcalib2026}. Its benchmark supplies composite memory blocks and annotates intended influence at the atomic level. Its training method separates overuse and underuse signals and redistributes credit using counterfactual localization. These results evaluate use of supplied memory, not the quality of a complete storage-and-retrieval system. The distinction limits the claim while making the post-retrieval problem experimentally visible.

Influence should follow role and authority. A product page can provide a claimed feature, a paper can provide an author-reported result, and a user message can provide a task constraint. Packing all three as equally authoritative text invites category errors. The model-visible representation should preserve source type and scope. It should also mark deductions separately from observations so that later summarization does not promote a proposed explanation into cited evidence.

Attribution requires a mapping from an output claim or action to supporting evidence. For a text answer, this can be a claim-span-to-source relation. For tool execution, it can include the source of each argument and the authority permitting its effect. A citation to a long document is not sufficient when the cited passage does not entail the claim. Source availability, source relevance, and claim support are different checks.

Conflicting evidence should remain visible until resolved. A packer that chooses only the highest-scoring passage can erase a known disagreement and create false certainty. A proposed packing objective reserves space for support, relevant contradiction, and unresolved uncertainty. This can reduce the number of distinct positive passages that fit, so its value should be evaluated through decision correctness and calibration rather than raw recall. The best evidence set can contain a passage that argues against the tentative answer.

The implementation should log both retrieved candidates and the final injected set. Record which records were excluded for budget, access, duplication, staleness, or low relevance. This supports targeted failure analysis and prevents a system from claiming that a fact was ``retrieved'' merely because it appeared in an internal candidate list. The final conditional input is the evidence boundary that matters for the policy.

